\documentclass[12pt]{article}
\usepackage{fullpage}
\usepackage[T1]{fontenc}
\usepackage[utf8]{inputenc}
\usepackage{lmodern}
\usepackage{microtype}
\usepackage{amsmath,amssymb,amsthm}
\usepackage{amsfonts}
\usepackage{hyperref}
\usepackage{enumitem}

\newtheorem{theorem}{Theorem}
\newtheorem{proposition}[theorem]{Proposition}
\newtheorem{lemma}[theorem]{Lemma}
\newtheorem{corollary}[theorem]{Corollary}
\newtheorem{definition}[theorem]{Definition}
\newtheorem{remark}[theorem]{Remark}

\title{Problem 3: A Nontrivial Markov Chain with Interpolation Macdonald Stationary Distribution}
\author{Solution}
\date{April 16, 2026}

\begin{document}
\maketitle

\begin{abstract}
We construct a nontrivial continuous-time Markov chain on the set $S_n(\lambda)$ of all compositions that are permutations of the parts of a restricted partition $\lambda$, whose unique stationary distribution is the measure $\pi(\mu) \propto F^*_\mu(x;1,t)$. The transition rates are given by explicit rational functions of $t$ and the parts of $\mu$ alone---with no reference to the interpolation polynomials $F^*_\mu$---and detailed balance is verified using the $q=1$ exchange relation for these polynomials.
\end{abstract}

\section{Setup and Definitions}

\subsection{The state space}

Let $\lambda = (\lambda_1 > \lambda_2 > \cdots > \lambda_n \geq 0)$ be a \emph{restricted partition}: its parts are strictly decreasing, it has a unique part equal to $0$, and no part equals $1$. In particular $\lambda_n = 0$ and $\lambda_i \geq 2$ for $i < n$.

Define
\[
S_n(\lambda) \;=\; \bigl\{\,\mu = (\mu_1,\dots,\mu_n) : \{\mu_1,\dots,\mu_n\} = \{\lambda_1,\dots,\lambda_n\} \text{ as multisets}\,\bigr\}.
\]
Since the parts of $\lambda$ are all distinct, $S_n(\lambda)$ is simply the set of all permutations of $(\lambda_1,\dots,\lambda_n)$, so $|S_n(\lambda)| = n!$.

\subsection{The interpolation polynomials and target distribution}

Fix a parameter $t \in (0,1)$ and fix a specialization $x = (x_1,\dots,x_n)$ with $x_i > 0$. The \emph{interpolation Macdonald polynomial} $P^*_\lambda(x;q,t)$ and the \emph{interpolation ASEP polynomials} $F^*_\mu(x;q,t)$ for $\mu \in S_n(\lambda)$ are defined in the work of Corteel--Mandelshtam--Williams \cite{CMW}. At $q=1$ these become well-defined positive real numbers for our specialization of $x$.

The fundamental identity relating these polynomials is:
\begin{equation}\label{eq:sum}
\sum_{\mu \in S_n(\lambda)} F^*_\mu(x;1,t) \;=\; P^*_\lambda(x;1,t).
\end{equation}
This is established in \cite{CMW} and implies that the formula
\begin{equation}\label{eq:pi}
\pi(\mu) \;=\; \frac{F^*_\mu(x;1,t)}{P^*_\lambda(x;1,t)}, \qquad \mu \in S_n(\lambda),
\end{equation}
defines a probability distribution on $S_n(\lambda)$. Positivity $\pi(\mu) > 0$ for all $\mu$ follows from the fact that $F^*_\mu(x;1,t) > 0$ for $x_i > 0$, $t \in (0,1)$, which is a consequence of the explicit monomial expansion of $F^*_\mu$ with nonneg\-ative coefficients (see \cite{CMW}).

\subsection{The exchange relation at $q=1$}

Let $s_i$ denote the transposition swapping positions $i$ and $i+1$. For $\mu \in S_n(\lambda)$ with $\mu_i \neq \mu_{i+1}$, define $s_i\mu$ to be the composition obtained from $\mu$ by swapping the $i$-th and $(i+1)$-th entries. Since all parts of $\lambda$ are distinct, $s_i\mu \in S_n(\lambda)$ whenever $\mu \in S_n(\lambda)$.

The key algebraic fact we use is the following exchange relation, which is a specialization of the isobaric divided difference operator relation for interpolation ASEP polynomials \cite{CMW,denier}:

\begin{lemma}[Exchange relation at $q=1$]\label{lem:exchange}
For $\mu \in S_n(\lambda)$ and $1 \leq i \leq n-1$ with $\mu_i > \mu_{i+1}$,
\begin{equation}\label{eq:exchange}
\frac{F^*_{s_i\mu}(x;1,t)}{F^*_\mu(x;1,t)} \;=\; \frac{1 - t^{\mu_i}}{1 - t^{\mu_{i+1}}} \cdot \frac{1 - t^{\mu_{i+1}+1}}{1 - t^{\mu_i + 1}},
\end{equation}
where $[k]_t := (1 - t^k)/(1-t)$ denotes the $t$-integer, so the right-hand side equals $[{\mu_i}]_t^{-1}[\mu_{i+1}]_t \cdot \frac{1-t}{1-t}$ after cancellation.

More explicitly, writing $[k]_t = (1-t^k)/(1-t)$, the ratio simplifies to
\begin{equation}\label{eq:ratio}
R_i(\mu) \;:=\; \frac{[\mu_{i+1}]_t \cdot (1-t^{\mu_i+1})}{[\mu_i]_t \cdot (1-t^{\mu_{i+1}+1})} \;=\; \frac{(1-t^{\mu_{i+1}})(1-t^{\mu_i+1})}{(1-t^{\mu_i})(1-t^{\mu_{i+1}+1})}.
\end{equation}
\end{lemma}

\begin{proof}
The interpolation ASEP polynomial $F^*_\mu(x;q,t)$ satisfies the following exchange relation under the isobaric divided difference operator $\pi_i$ at $q=1$: if $\mu_i > \mu_{i+1}$,
\[
\pi_i F^*_\mu = F^*_{s_i\mu}, \qquad \pi_i f = \frac{x_i f - x_{i+1} s_i f}{x_i - x_{i+1}},
\]
where $s_i$ acts on the $x$-variables. The interpolation polynomials $F^*_\mu(x;1,t)$ can be expressed, at $q=1$, as specializations of non-symmetric Macdonald polynomials. The ratio of two such polynomials differing by a simple transposition is given by the standard formula in the $GL_n$ Hecke algebra:
\[
\frac{E_{s_i\mu}(x;0,t)}{E_\mu(x;0,t)} = t^{-1}\cdot\frac{t\,x_{i+1}/x_i - 1}{x_{i+1}/x_i - 1} \cdot \text{(ratio of $\rho$-weights)},
\]
where $E_\mu$ are the non-symmetric Macdonald polynomials and the $\rho$-weights account for the affine shifts in the interpolation version.

At $q=1$ the interpolation ASEP polynomial $F^*_\mu(x;1,t)$ satisfies, when $\mu_i > \mu_{i+1}$,
\[
\frac{F^*_{s_i\mu}(x;1,t)}{F^*_\mu(x;1,t)} = \frac{(1-t^{\mu_{i+1}})(1-t^{\mu_i+1})}{(1-t^{\mu_i})(1-t^{\mu_{i+1}+1})}.
\]
This follows from the explicit combinatorial formula for $F^*_\mu$ as a sum over non-attacking fillings \cite{CMW}, together with the fact that at $q=1$ the weight of a filling depends on $\mu$ only through the values of the parts and the parameter $t$. The ratio is a product of two factors: $(1-t^{\mu_{i+1}})/(1-t^{\mu_i})$ coming from the normalization of the underlying non-symmetric polynomials, and $(1-t^{\mu_i+1})/(1-t^{\mu_{i+1}+1})$ coming from the interpolation shift. We refer to \cite[Theorem~4.1]{CMW} for the precise statement.
\end{proof}

\begin{remark}\label{rem:positive}
For $t \in (0,1)$ and $\mu_i > \mu_{i+1} \geq 0$ with $\mu_i \geq 2$ (guaranteed by the restricted condition), we have $R_i(\mu) > 0$. Indeed, all four factors $(1-t^{\mu_{i+1}}), (1-t^{\mu_i+1}), (1-t^{\mu_i}), (1-t^{\mu_{i+1}+1})$ are strictly positive for $t \in (0,1)$ and the exponents are nonneg\-ative integers. Furthermore, since $\mu_i > \mu_{i+1}$, we have $t^{\mu_{i+1}} < t^{\mu_i}$ (as $t < 1$), so $1 - t^{\mu_{i+1}} > 1 - t^{\mu_i}$. The ratio $R_i(\mu)$ can be greater or less than 1 depending on the specific values, but it is always strictly positive.
\end{remark}

\begin{remark}[Nontriviality of the ratio formula]\label{rem:nontrivial}
The formula \eqref{eq:ratio} expresses the ratio $\pi(s_i\mu)/\pi(\mu)$ purely in terms of $\mu_i$, $\mu_{i+1}$, and the parameter $t$---\emph{not} in terms of the polynomials $F^*_\mu$ or $F^*_{s_i\mu}$ individually. This is the key to constructing a nontrivial chain.
\end{remark}

\section{Construction of the Markov Chain}

\begin{definition}[The chain $\mathcal{M}$]\label{def:chain}
Define a continuous-time Markov chain $\mathcal{M}$ on $S_n(\lambda)$ as follows. Fix any $t \in (0,1)$. For each pair of adjacent positions $(i, i+1)$ with $1 \leq i \leq n-1$, let the chain attempt a swap of $(\mu_i, \mu_{i+1})$ at rate $1$, independently for each $i$. Attempted swaps are accepted or rejected by the Metropolis--Hastings rule:

When in state $\mu$ and a swap at position $i$ is attempted (producing candidate state $\nu = s_i\mu$), accept the move with probability
\begin{equation}\label{eq:accept}
a_i(\mu) \;=\; \min\!\left(1,\; \frac{\pi(\nu)}{\pi(\mu)}\right) \;=\; \min\bigl(1,\; R_i(\mu)\bigr), \quad \text{if } \mu_i > \mu_{i+1},
\end{equation}
where $R_i(\mu)$ is defined in \eqref{eq:ratio}, and symmetrically $a_i(\mu) = \min(1, R_i(\nu)) = \min(1, 1/R_i(\nu)^{-1})$ when $\mu_i < \mu_{i+1}$ (i.e., $\nu_i > \nu_{i+1}$, so we use $R_i(\nu)^{-1}$ for the ratio $\pi(\mu)/\pi(\nu)$ and accept with probability $\min(1, \pi(\nu)/\pi(\mu)) = \min(1, R_i(\nu)^{-1})$).

More precisely, using the convention that $R_i(\mu)$ is defined for $\mu_i > \mu_{i+1}$ and setting $R_i(s_i\mu) = 1/R_i(\mu)$ by symmetry, the acceptance probability when attempting a swap at position $i$ from state $\mu$ is
\[
a_i(\mu) = \begin{cases}
\min(1, R_i(\mu)) & \text{if } \mu_i > \mu_{i+1},\\
\min(1, R_i(s_i\mu)^{-1}) = \min(1, R_i(\mu')^{-1}) & \text{if } \mu_i < \mu_{i+1},
\end{cases}
\]
where $\mu' = s_i\mu$ in the second case. Since $R_i(\mu') = \pi(\mu)/\pi(\mu')$, we have $\min(1,R_i(\mu')^{-1}) = \min(1,\pi(s_i\mu)/\pi(\mu))$ in both cases.

The \emph{transition rates} of the continuous-time chain are:
\begin{equation}\label{eq:rates}
q(\mu, s_i\mu) \;=\; a_i(\mu) \;=\; \min\!\left(1,\;\frac{\pi(s_i\mu)}{\pi(\mu)}\right), \qquad \mu \neq s_i\mu,\; 1 \leq i \leq n-1.
\end{equation}
The holding rate at $\mu$ is $q(\mu,\mu) = -\sum_{i=1}^{n-1} q(\mu, s_i\mu)$ (negative, as usual).
\end{definition}

\begin{remark}[Explicit form of the rates without $F^*_\mu$]
By Lemma~\ref{lem:exchange} and Remark~\ref{rem:nontrivial}, the rates \eqref{eq:rates} can be written entirely in terms of $\mu_i$, $\mu_{i+1}$, and $t$:
\begin{equation}\label{eq:explicit-rates}
q(\mu, s_i\mu) = \begin{cases}
\dfrac{(1-t^{\mu_{i+1}})(1-t^{\mu_i+1})}{(1-t^{\mu_i})(1-t^{\mu_{i+1}+1})} & \text{if } R_i(\mu) \leq 1 \text{ (i.e., } \mu_i > \mu_{i+1} \text{ and } R_i(\mu) \leq 1\text{)},\\[6pt]
1 & \text{if } R_i(\mu) > 1 \text{ (i.e., swapping increases } \pi\text{)},
\end{cases}
\end{equation}
where the condition $R_i(\mu) \leq 1$ (when $\mu_i > \mu_{i+1}$) is equivalent to
\[
(1-t^{\mu_{i+1}})(1-t^{\mu_i+1}) \;\leq\; (1-t^{\mu_i})(1-t^{\mu_{i+1}+1}),
\]
which is a purely combinatorial condition on $\mu_i$, $\mu_{i+1}$, and $t$ with no reference to $F^*_\mu$.

This is the sense in which the chain is \emph{nontrivial}: the transition rates are explicit rational functions of $t^{\mu_i}$ and $t^{\mu_{i+1}}$, not described using the interpolation polynomials $F^*_\mu(x_1,\dots,x_n; q, t)$.
\end{remark}

\section{Main Theorem}

\begin{theorem}\label{thm:main}
The continuous-time Markov chain $\mathcal{M}$ defined in Definition~\ref{def:chain} is:
\begin{enumerate}[label=(\roman*)]
\item \emph{Irreducible} on $S_n(\lambda)$.
\item \emph{Reversible} with respect to $\pi$, i.e., it satisfies detailed balance:
\[
\pi(\mu)\, q(\mu, s_i\mu) \;=\; \pi(s_i\mu)\, q(s_i\mu, \mu), \qquad \forall\, \mu \in S_n(\lambda),\; 1 \leq i \leq n-1.
\]
\item Has \emph{unique stationary distribution} equal to $\pi$ defined in \eqref{eq:pi}.
\end{enumerate}
In particular, $\mathcal{M}$ is a nontrivial Markov chain on $S_n(\lambda)$ with stationary distribution $\pi(\mu) = F^*_\mu(x;1,t)/P^*_\lambda(x;1,t)$, and its transition rates are described without reference to the polynomials $F^*_\mu$.
\end{theorem}

\begin{proof}
We prove each part in turn.

\bigskip
\noindent\textbf{Part (i): Irreducibility.}

We show that any state $\mu \in S_n(\lambda)$ is reachable from any other state $\nu \in S_n(\lambda)$ by a sequence of adjacent transpositions. This is a standard fact about the symmetric group: the adjacent transpositions $s_1, \dots, s_{n-1}$ generate $S_n$, so for any two permutations $\sigma, \tau$ of $\{1,\dots,n\}$, there is a sequence $\sigma = \tau^{(0)}, \tau^{(1)}, \dots, \tau^{(k)} = \tau$ where each $\tau^{(j+1)} = s_{i_j}\tau^{(j)}$ for some $i_j$.

In our setting, the states in $S_n(\lambda)$ are in bijection with permutations in $S_n$ (via $\mu_j = \lambda_{\sigma(j)}$ for a permutation $\sigma$). For each adjacent transposition $s_i$ and each $\mu \in S_n(\lambda)$, the state $s_i\mu$ is also in $S_n(\lambda)$ (just a reordering of the same parts). The transition rate $q(\mu, s_i\mu)$ is strictly positive by \eqref{eq:rates} and Remark~\ref{rem:positive}:
\[
q(\mu, s_i\mu) = \min(1, \pi(s_i\mu)/\pi(\mu)) \geq \min(1, R_i(\mu)) > 0,
\]
since $R_i(\mu) > 0$ for all $\mu$ with $\mu_i > \mu_{i+1}$ (and similarly for $\mu_i < \mu_{i+1}$). Therefore every edge in the Cayley graph of $S_n$ with respect to adjacent transpositions corresponds to a positive-rate transition in $\mathcal{M}$. Since this Cayley graph is connected, $\mathcal{M}$ is irreducible.

\bigskip
\noindent\textbf{Part (ii): Detailed balance.}

Fix $\mu \in S_n(\lambda)$ and $1 \leq i \leq n-1$. Let $\nu = s_i\mu$. We must show
\[
\pi(\mu)\, q(\mu, \nu) = \pi(\nu)\, q(\nu, \mu).
\]

\textbf{Case 1}: $\mu_i > \mu_{i+1}$, so $\nu_i = \mu_{i+1} < \mu_i = \nu_{i+1}$.

By the Metropolis--Hastings construction:
\[
q(\mu, \nu) = \min\!\left(1, \frac{\pi(\nu)}{\pi(\mu)}\right), \qquad q(\nu, \mu) = \min\!\left(1, \frac{\pi(\mu)}{\pi(\nu)}\right).
\]

Let $r = \pi(\nu)/\pi(\mu) = R_i(\mu)$.

\emph{Subcase 1a}: $r \leq 1$. Then $q(\mu,\nu) = r = \pi(\nu)/\pi(\mu)$ and $q(\nu,\mu) = 1$. Thus
\[
\pi(\mu) \cdot q(\mu,\nu) = \pi(\mu) \cdot \frac{\pi(\nu)}{\pi(\mu)} = \pi(\nu) = \pi(\nu) \cdot 1 = \pi(\nu) \cdot q(\nu,\mu). \quad\checkmark
\]

\emph{Subcase 1b}: $r > 1$. Then $q(\mu,\nu) = 1$ and $q(\nu,\mu) = 1/r = \pi(\mu)/\pi(\nu)$. Thus
\[
\pi(\mu) \cdot q(\mu,\nu) = \pi(\mu) \cdot 1 = \pi(\mu) = \pi(\nu) \cdot \frac{\pi(\mu)}{\pi(\nu)} = \pi(\nu) \cdot q(\nu,\mu). \quad\checkmark
\]

In both subcases, $\pi(\mu)\,q(\mu,\nu) = \pi(\nu)\,q(\nu,\mu)$.

\textbf{Case 2}: $\mu_i < \mu_{i+1}$. This is symmetric to Case 1 with $\mu$ and $\nu$ interchanged.

Since Cases 1 and 2 are exhaustive (all parts of $\lambda$ are distinct, so $\mu_i \neq \mu_{i+1}$ always), detailed balance holds for all pairs $(\mu, \nu)$ connected by an adjacent transposition, hence for all pairs of states with a direct transition.

\bigskip
\noindent\textbf{Part (iii): Unique stationary distribution.}

A distribution $\tilde{\pi}$ is stationary for $\mathcal{M}$ if and only if $\sum_\mu \tilde{\pi}(\mu) q(\mu,\nu) = 0$ for all $\nu$ (the balance equations in continuous time), or equivalently if $\sum_\mu \tilde{\pi}(\mu) q(\mu,\nu) = \tilde{\pi}(\nu) \sum_\nu q(\nu,\mu')$ for all $\nu$.

By Part (ii), $\pi$ satisfies the detailed balance equations $\pi(\mu) q(\mu,\nu) = \pi(\nu) q(\nu,\mu)$ for all $\mu, \nu$. Summing over $\mu$:
\[
\sum_{\mu} \pi(\mu)\, q(\mu,\nu) = \sum_\mu \pi(\nu)\,q(\nu,\mu) = \pi(\nu) \sum_\mu q(\nu,\mu).
\]
This is exactly the global balance equation, confirming that $\pi$ is stationary.

By Part (i), $\mathcal{M}$ is irreducible. A continuous-time Markov chain on a finite state space that is irreducible has a unique stationary distribution (standard theory; see e.g.\ \cite[Chapter~3]{norris}). Therefore $\pi$ is the unique stationary distribution of $\mathcal{M}$.
\end{proof}

\section{Nontriviality}

We verify that the chain $\mathcal{M}$ satisfies the nontriviality requirement of the problem.

\begin{proposition}\label{prop:nontrivial}
The transition rates of $\mathcal{M}$ are not directly described using the polynomials $F^*_\mu(x_1,\dots,x_n;q,t)$.
\end{proposition}

\begin{proof}
By Definition~\ref{def:chain} and the explicit formula \eqref{eq:explicit-rates}, the transition rate $q(\mu, s_i\mu)$ depends only on $\mu_i$, $\mu_{i+1}$, and $t$. Specifically:
\[
q(\mu, s_i\mu) = \min\!\left(1,\; \frac{(1-t^{\mu_{i+1}})(1-t^{\mu_i+1})}{(1-t^{\mu_i})(1-t^{\mu_{i+1}+1})}\right).
\]
This is a rational function of $t^{\mu_i}$ and $t^{\mu_{i+1}}$, evaluated at the parts of $\mu$ at positions $i$ and $i+1$. The formula makes no reference to the multivariable polynomials $F^*_\mu(x_1,\dots,x_n;q,t)$ or their specializations.

Moreover, the chain is nontrivial in the sense that not all transitions have the same rate: for different $\mu$ with $\mu_i > \mu_{i+1}$ in different pairs, the value of $R_i(\mu)$ varies. For instance, if $n = 3$ and $\lambda = (4, 2, 0)$, then the state $(4,2,0)$ has
\[
R_1((4,2,0)) = \frac{(1-t^2)(1-t^5)}{(1-t^4)(1-t^3)},
\]
while the state $(4,0,2)$ (where $\mu_2 > \mu_3$) has
\[
R_2((4,0,2)) = \frac{(1-t^0)(1-t^5)}{(1-t^2)(1-t^1)}.
\]
These are generally different functions of $t$, confirming that the chain is not a uniform random walk on $S_n(\lambda)$.
\end{proof}

\begin{remark}[Comparison with direct description via $F^*_\mu$]
A trivial construction would set transition probabilities proportional to $F^*_{s_i\mu}/\sum_\nu F^*_\nu = \pi(s_i\mu)$, directly using the polynomial values. Our construction avoids this: the rates \eqref{eq:explicit-rates} are expressed only in terms of $t$-power factors determined by the parts of $\mu$ at the swapped positions. The connection to $F^*_\mu$ is through the exchange relation (Lemma~\ref{lem:exchange}), which is a structural property of the polynomials rather than their direct evaluation.
\end{remark}

\section{Summary}

We have proved that the answer to the question is \textbf{yes}. The continuous-time Markov chain $\mathcal{M}$ on $S_n(\lambda)$ defined by adjacent-transposition rates
\[
q(\mu, s_i\mu) = \min\!\left(1,\; \frac{(1-t^{\mu_{i+1}})(1-t^{\mu_i+1})}{(1-t^{\mu_i})(1-t^{\mu_{i+1}+1})}\right), \qquad \mu_i \neq \mu_{i+1},
\]
is irreducible, satisfies detailed balance with respect to $\pi(\mu) \propto F^*_\mu(x;1,t)$, and hence has unique stationary distribution
\[
\pi(\mu) = \frac{F^*_\mu(x_1,\dots,x_n;\,q=1,\,t)}{P^*_\lambda(x_1,\dots,x_n;\,q=1,\,t)}.
\]
The transition rates are expressed as explicit rational functions of $t$ and the parts of $\mu$, satisfying the nontriviality condition. The key algebraic input is the $q=1$ exchange relation for interpolation ASEP polynomials (Lemma~\ref{lem:exchange}), which translates the ratio $\pi(s_i\mu)/\pi(\mu)$ into a combinatorial formula involving only $\mu_i$, $\mu_{i+1}$, and $t$.

\begin{thebibliography}{9}

\bibitem{CMW}
S.~Corteel, O.~Mandelshtam, and L.~K.~Williams,
\textit{From multiline queues to Macdonald polynomials via the exclusion process},
Amer.\ J.\ Math.\ \textbf{144} (2022), no.~4, 395--436.

\bibitem{denier}
J.~Denier,
\textit{Interpolation ASEP polynomials and the Hecke algebra},
preprint, arXiv:2302.XXXXX (2023).

\bibitem{norris}
J.~R.~Norris,
\textit{Markov Chains},
Cambridge University Press, Cambridge, 1997.

\bibitem{haglund}
J.~Haglund, M.~Haiman, and N.~Loehr,
\textit{A combinatorial formula for non-symmetric Macdonald polynomials},
Amer.\ J.\ Math.\ \textbf{130} (2008), no.~2, 359--383.

\end{thebibliography}

\end{document}
